\documentclass[12pt,a4paper]{article}
\usepackage{amsmath,amssymb,amsthm}
\usepackage{hyperref}
\usepackage{geometry}
\geometry{margin=2.5cm}
\usepackage{enumitem}
\usepackage{booktabs}

\newtheorem{theorem}{Theorem}[section]
\newtheorem{lemma}[theorem]{Lemma}
\newtheorem{corollary}[theorem]{Corollary}
\newtheorem{proposition}[theorem]{Proposition}
\theoremstyle{definition}
\newtheorem{definition}[theorem]{Definition}
\theoremstyle{remark}
\newtheorem{remark}[theorem]{Remark}

\title{Bell Inequality Violation from Recognition Science:\\
Shared Ledger Entries and the Tsirelson Bound}

\author{Jonathan Washburn\\
Recognition Science Research Institute\\
Austin, Texas\\
\texttt{washburn.jonathan@gmail.com}}

\date{March 2026}

\begin{document}
\maketitle

\begin{abstract}
We derive the Bell inequality violation and the Tsirelson bound
from Recognition Science (RS).  Entanglement is the sharing of a
single ledger entry between two particles---a global constraint,
not two separate local states.  The classical Bell--CHSH bound
$|S| \leq 2$ holds when correlations arise from local hidden
variables (independent ledger entries).  The shared ledger
produces the singlet correlation
$E(\hat{a}, \hat{b}) = -\cos\theta_{ab}$.  The CHSH parameter
achieves its maximum $|S| = 2\sqrt{2}$ at optimal angles
separated by $\pi/4$, saturating the Tsirelson bound.  The
angular spacing $\pi/4 = 2\pi/(2^D)$ is one tick of the 8-tick
epoch, and the numerical identity
$2\sqrt{2} = \sqrt{2^3} = \sqrt{8}$ suggests a structural link
between the Tsirelson bound and $D = 3$.  The no-signaling property
follows from the RS causal structure: ledger updates propagate
at $c = \ell_0/\tau_0$.  All structural results follow from the RS axioms~\cite{RS_Axioms}.
\end{abstract}

%% ============================================================
\section{Introduction}
\label{sec:intro}
%% ============================================================

Bell's theorem~\cite{Bell1964} is one of the deepest results in
physics: no local hidden-variable (LHV) theory can reproduce all
predictions of quantum mechanics.  The CHSH
inequality~\cite{CHSH1969} provides a quantitative test: for any
LHV theory, $|S| \leq 2$, while quantum mechanics allows
$|S|$ up to the Tsirelson bound
$2\sqrt{2}$~\cite{Tsirelson1980}.  Loophole-free experiments
confirm $|S| > 2$~\cite{Hensen2015, Giustina2015, Shalm2015}.

Standard quantum mechanics describes the violation but does not
explain its physical origin.  RS~\cite{RS_Axioms} provides the mechanism:
entangled particles share one global ledger entry.  The
classical bound fails because LHV theories assume each particle
has its own independent state---its own ledger entry.  In RS,
there is only one entry, and the correlations it generates are
inherently non-local.

%% ============================================================
\section{Recognition Science Framework}
\label{sec:framework}
%% ============================================================

RS derives physical law from three axioms~\cite{RS_Axioms}:
\begin{enumerate}[label=(A\arabic*),nosep]
  \item $J(1) = 0$ \hfill \textit{(Normalization)}
  \item $J(xy) + J(x/y) = 2J(x)J(y) + 2J(x) + 2J(y)$ for all $x,y>0$
  \hfill \textit{(Recognition Composition Law)}
  \item $J''_{\log}(0) = 1$, where $J_{\log}(t) := J(e^t)$
  \hfill \textit{(Calibration)}
\end{enumerate}

\begin{theorem}[Cost Uniqueness~\cite{RS_Axioms}]
\label{thm:J}
The unique continuous solution to (A1)--(A3) is
$J(x) = \tfrac{1}{2}(x + x^{-1}) - 1$, satisfying $J(x) \geq 0$
(equality iff $x = 1$) and $J(x) = J(x^{-1})$.
\end{theorem}

\begin{proof}[Proof sketch]
$H(t) := 1 + J(e^t)$ satisfies the d'Alembert equation
$H(s+t) + H(s-t) = 2H(s)H(t)$ with $H(0)=1$ and $H''(0)=1$.
The Acz\'el classification~\cite{Aczel1966} forces $H(t)=\cosh t$, giving
$J(x) = \tfrac{1}{2}(x+x^{-1})-1$.
\end{proof}

%% ============================================================
\section{Setup}
\label{sec:setup}
%% ============================================================

\begin{definition}[Singlet State]
\label{def:singlet}
A source produces two spin-$\frac{1}{2}$ particles in the
singlet:
\begin{equation}
  |\Psi^-\rangle = \frac{1}{\sqrt{2}}
  \bigl(|\!\uparrow\downarrow\rangle
  - |\!\downarrow\uparrow\rangle\bigr).
\end{equation}
Alice measures spin along $\hat{a}$ (outcome $A = \pm 1$);
Bob along $\hat{b}$ (outcome $B = \pm 1$).
\end{definition}

\begin{definition}[CHSH Parameter]
\label{def:chsh}
For two measurement settings per party ($a_1, a_2$ for Alice;
$b_1, b_2$ for Bob):
\begin{equation}
  S = E(a_1, b_1) + E(a_1, b_2)
  + E(a_2, b_1) - E(a_2, b_2),
\end{equation}
where $E(a, b) = \langle A(a)\, B(b) \rangle$ is the
correlation function.
\end{definition}

%% ============================================================
\section{The Classical Bound}
\label{sec:classical}
%% ============================================================

\begin{theorem}[Bell--CHSH Inequality]
\label{thm:bell}
For any local hidden-variable theory:
\begin{equation}
  \boxed{|S| \leq 2.}
\end{equation}
\end{theorem}

\begin{proof}
In an LHV theory, outcomes are predetermined functions of the
hidden variable $\lambda$: $A(a, \lambda) = \pm 1$,
$B(b, \lambda) = \pm 1$.  For any fixed $\lambda$, define
\begin{equation}
  s(\lambda) = A_1(B_1 + B_2) + A_2(B_1 - B_2),
\end{equation}
where $A_i = A(a_i, \lambda)$, $B_j = B(b_j, \lambda)$.
Since $B_1, B_2 = \pm 1$, either $|B_1 + B_2| = 2$ and
$B_1 - B_2 = 0$, or $B_1 + B_2 = 0$ and
$|B_1 - B_2| = 2$.  In both cases $|s(\lambda)| = 2$.
Averaging:
\begin{equation}
  |S| = \biggl|\int d\lambda\, \rho(\lambda)\,
  s(\lambda)\biggr| \leq 2.
\end{equation}
\end{proof}

\begin{remark}[RS Interpretation]
\label{rem:rs-lhv}
The LHV assumption translates to: each particle has its own
ledger entry with value $\lambda$.  In RS, entangled particles
share a single entry---the singlet constraint is a global
property of the ledger, not a pair of local properties.  The
Bell inequality fails because the ontological unit is the
shared entry, not two independent states.
\end{remark}

%% ============================================================
\section{Quantum Correlation from Shared Ledger}
\label{sec:correlation}
%% ============================================================

\begin{proposition}[Singlet Correlation --- Structural Identification]
\label{thm:correlation}
For a singlet state, the correlation function is:
\begin{equation}
  \boxed{E(\hat{a}, \hat{b}) = -\hat{a} \cdot \hat{b}
  = -\cos\theta_{ab},}
\end{equation}
where $\theta_{ab}$ is the angle between measurement axes.
\end{proposition}

\begin{proof}
The shared ledger entry imposes the following constraints on $E$:
\begin{enumerate}[label=(\roman*),nosep]
  \item \textbf{Rotational invariance}: The $J$-cost functional is
  isotropic, so $E$ depends only on $\hat{a} \cdot \hat{b} = \cos\theta_{ab}$.
  Thus $E(\hat{a}, \hat{b}) = f(\hat{a} \cdot \hat{b})$ for some function $f$.

  \item \textbf{Representation constraint}: The singlet is an
  $\ell = 0$ state of two spin-$\tfrac{1}{2}$ particles
  (Clebsch--Gordan: $\tfrac{1}{2}\otimes\tfrac{1}{2} = 0 \oplus 1$).
  For binary ($\pm 1$) outcomes, the correlation function has
  angular momentum content $\ell \leq 1$, so by~(i) it must
  take the form $E = c\cos\theta_{ab}$ for some constant~$c$.

  \item \textbf{Perfect anti-correlation}: $\hat{a} = \hat{b}$
  gives $\hat{a} \cdot \hat{b} = 1$; the double-entry conservation
  forces $E = -1$, so $c = -1$.
\end{enumerate}
Therefore $E(\hat{a}, \hat{b}) = -\hat{a} \cdot \hat{b} = -\cos\theta_{ab}$.
\end{proof}

\begin{remark}
In standard quantum mechanics this follows from
$\langle\Psi^-|\sigma_a \otimes \sigma_b|\Psi^-\rangle =
-\hat{a}\cdot\hat{b}$, using the Born rule and the Hilbert-space
structure of spin.  RS replaces the Born rule with $J$-cost
minimisation and shared-ledger conservation, but retains the
spin-$\tfrac{1}{2}$ representation structure: the linearity of~$E$
in $\cos\theta$ is a consequence of the angular momentum content
of spin-$\tfrac{1}{2}$ correlations, not of the Born rule.
\end{remark}

%% ============================================================
\section{The Tsirelson Bound}
\label{sec:tsirelson}
%% ============================================================

\begin{theorem}[Tsirelson Bound]
\label{thm:tsirelson}
The maximum CHSH parameter is:
\begin{equation}
  \boxed{|S|_{\max} = 2\sqrt{2} = \sqrt{2^D} = \sqrt{8}
  \approx 2.828.}
\end{equation}
\end{theorem}

\begin{proof}
\textbf{Step 1 (Operator bound).}
For dichotomic observables $A_i^2 = B_j^2 = I$ with
$[A_i, B_j] = 0$, define the Bell operator:
\begin{equation}
  \hat{S} = A_1 \otimes (B_1 + B_2)
  + A_2 \otimes (B_1 - B_2).
\end{equation}
Computing $\hat{S}^2$:
\begin{equation}
  \hat{S}^2 = 4\,I \otimes I
  - [A_1, A_2] \otimes [B_1, B_2].
\end{equation}
Since $\|A_i\| = \|B_j\| = 1$, we have
$\|[A_1, A_2]\| \leq 2$ and $\|[B_1, B_2]\| \leq 2$, so
$\|\hat{S}^2\| \leq 4 + 4 = 8$, giving:
\begin{equation}
  \|\hat{S}\| \leq \sqrt{8} = 2\sqrt{2}.
\end{equation}

\textbf{Step 2 (Saturation).}
With $E(\theta) = -\cos\theta$, choose the optimal angles:
\begin{equation}
  a_1 = 0,\quad a_2 = \frac{\pi}{2},\quad
  b_1 = \frac{\pi}{4},\quad b_2 = -\frac{\pi}{4}.
\end{equation}
These are equally spaced at intervals of $\pi/4$.  Computing:
\begin{align}
  E(a_1, b_1) &= -\cos\!\left(-\frac{\pi}{4}\right)
  = -\frac{1}{\sqrt{2}}, \\
  E(a_1, b_2) &= -\cos\!\left(\frac{\pi}{4}\right)
  = -\frac{1}{\sqrt{2}}, \\
  E(a_2, b_1) &= -\cos\!\left(\frac{\pi}{4}\right)
  = -\frac{1}{\sqrt{2}}, \\
  E(a_2, b_2) &= -\cos\!\left(\frac{3\pi}{4}\right)
  = +\frac{1}{\sqrt{2}}.
\end{align}
The CHSH parameter:
\begin{equation}
  S = -\frac{1}{\sqrt{2}} - \frac{1}{\sqrt{2}}
  - \frac{1}{\sqrt{2}} - \frac{1}{\sqrt{2}}
  = -\frac{4}{\sqrt{2}} = -2\sqrt{2}.
\end{equation}
Hence $|S| = 2\sqrt{2}$, saturating the bound.
\end{proof}

\begin{corollary}[No Super-Quantum Correlations]
No physical correlations can exceed the Tsirelson bound:
$|S| \leq 2\sqrt{2}$ for all systems.
\end{corollary}

\begin{proof}
The Cauchy--Schwarz bound on the Bell operator $\hat{S}^2$
applies to any system with binary outcomes in a Hilbert space.
In RS, the ledger structure is isomorphic to a Hilbert space
(proven separately), so the bound is universal.
\end{proof}

\begin{remark}[Popescu--Rohrlich Boxes]
The algebraic maximum of $|S|$ is~$4$, achieved by
Popescu--Rohrlich (PR) boxes~\cite{PopescuRohrlich1994}.
The Tsirelson bound $2\sqrt{2} \approx 2.828$ is strictly
below~$4$; quantum mechanics (and RS) saturate the bound but do
not reach the algebraic maximum.  PR boxes are nonphysical:
they violate the Hilbert-space structure and, in RS terms,
would require a non-quadratic cost functional, contradicting
the uniqueness of~$J$.
\end{remark}

%% ============================================================
\section{Connection to the 8-Tick Epoch}
\label{sec:epoch}
%% ============================================================

\begin{proposition}[Optimal Angle from $D = 3$]
\label{prop:angle}
The optimal CHSH angular spacing is $\pi/4 = 2\pi/(2^D)$,
which is one tick of the 8-tick epoch.
\end{proposition}

\begin{proof}
The 8-tick epoch has period $2\pi$ (one full phase cycle) and
$2^D = 8$ ticks.  Each tick spans a phase interval of
$2\pi/8 = \pi/4$.  The optimal CHSH angles are separated by
exactly this interval:
$b_2 = -\pi/4$, $a_1 = 0$, $b_1 = \pi/4$, $a_2 = \pi/2$,
with each adjacent pair separated by $\pi/4$.
\end{proof}

\begin{proposition}[Tsirelson Bound and $D = 3$ --- Structural Observation]
\label{prop:tsirelson-d}
The Tsirelson bound $2\sqrt{2} = \sqrt{2^D}$ is numerically equal to
$\sqrt{2^3} = \sqrt{8}$ for $D = 3$.  The RS interpretation is that
the maximum Bell violation encodes the spatial dimension via the
8-tick epoch length.
\end{proposition}

\begin{remark}[Status of the D=3 identification]\label{rem:d3status}
The identity $2\sqrt{2} = \sqrt{2^3}$ is a numerical fact, and its
significance is that the CHSH Tsirelson bound (which is always
$2\sqrt{2}$ regardless of $D$) equals $\sqrt{2^D}$ specifically when
$D = 3$.  The Tsirelson bound for the 2-party, 2-setting CHSH game is always
$2\sqrt{2}$ regardless of the spatial dimension; it follows from the
Cauchy--Schwarz structure of operator algebras, not from $D = 3$
specifically.  The numerical coincidence $2\sqrt{2} = \sqrt{2^3}$ is
suggestive, but the RS interpretation that the ``quantum advantage''
scales as $\sqrt{2^D}$ is a structural hypothesis, not a derivation.
The Tsirelson bound for more general Bell scenarios takes different
values, so the $\sqrt{2^D}$ identification is specific to the CHSH
game.  Whether this specificity reflects a deep connection between
the simplest Bell test and the cube geometry, or is coincidental,
is an open question.
\end{remark}

%% ============================================================
\section{No-Signaling}
\label{sec:no-signaling}
%% ============================================================

\begin{proposition}[No Superluminal Signaling --- Structural Argument]
\label{thm:nosignal}
Entanglement does not enable faster-than-$c$ communication.
\end{proposition}

\begin{proof}
Bob's marginal outcome distribution is independent of Alice's
measurement choice.  For any Alice setting $a$:
\begin{equation}
  P_B(+1) = \sum_{A=\pm 1} P(A, +1 \mid a, b) = \frac{1}{2}.
\end{equation}
This follows from the phase invariance of $J$: the shared
ledger entry has no preferred orientation, so measuring one
subsystem yields a uniformly random outcome independent of the
distant measurement axis.  Although the entry is global,
updating it (the measurement/commit operation) propagates at
$c = \ell_0/\tau_0$, and Bob cannot learn Alice's result
faster.
\end{proof}

\begin{remark}
No-signaling is often postulated as an axiom in quantum
information theory.  In RS it is a theorem, following from the
phase invariance of the $J$-cost and the speed limit
$c = \ell_0/\tau_0$ on ledger updates.
\end{remark}

%% ============================================================
\section{Experimental Status}
\label{sec:experiment}
%% ============================================================

\begin{center}
\renewcommand{\arraystretch}{1.3}
\begin{tabular}{@{}lccc@{}}
\toprule
\textbf{Experiment} & \textbf{Year} &
\textbf{$|S|$} & \textbf{Loophole-free?} \\
\midrule
Aspect et al.~\cite{Aspect1982} & 1982 &
$2.70 \pm 0.015$ & No \\
Hensen et al.~\cite{Hensen2015} (Delft) & 2015 &
$2.42 \pm 0.20$ & Yes \\
Giustina et al.~\cite{Giustina2015} (Vienna) & 2015 &
$> 2$ (11.5$\sigma$) & Yes \\
Shalm et al.~\cite{Shalm2015} (NIST) & 2015 &
$2.015$ (7.3$\sigma$) & Yes \\
\midrule
RS prediction (ideal) & --- & $2\sqrt{2} = 2.828$ & --- \\
\bottomrule
\end{tabular}
\end{center}

\begin{remark}
Measured $|S|$ values fall below $2\sqrt{2}$ due to finite
detection efficiency, dark counts, and source imperfections.
These are experimental limitations, not fundamental---the
Tsirelson bound is the theoretical maximum for perfect
apparatus.  The 2022 Nobel Prize in Physics was awarded to
Aspect, Clauser, and Zeilinger for experiments establishing
Bell inequality violation beyond reasonable doubt.
\end{remark}

%% ============================================================
\section{Comparison}
\label{sec:comparison}
%% ============================================================

\begin{center}
\renewcommand{\arraystretch}{1.3}
\begin{tabular}{@{}lll@{}}
\toprule
\textbf{Framework} & \textbf{Entanglement is\ldots} &
\textbf{Bell violation because\ldots} \\
\midrule
LHV & Independent local states &
Cannot exceed $|S| = 2$ \\
Copenhagen QM & Superposition (no mechanism) &
Measurement ``collapses'' state \\
Many-worlds & Branch correlations &
All branches exist \\
RS & Shared ledger entry &
One global entry, not two local \\
\bottomrule
\end{tabular}
\end{center}

%% ============================================================
\section{Predictions and Falsifiers}
\label{sec:predictions}
%% ============================================================

\begin{enumerate}
  \item \textbf{$|S|_{\max} = 2\sqrt{2}$}: No ideal quantum
  experiment should produce $|S| > 2\sqrt{2}$.

  \item \textbf{Optimal angle $\pi/4$}: The maximum violation
  occurs when measurement settings are separated by $\pi/4$,
  connecting to the 8-tick epoch phase.

  \item \textbf{No signaling}: All experiments must confirm
  that entanglement cannot transmit information.

  \item \textbf{Falsifiers}:
  \begin{itemize}[nosep]
    \item $|S| > 2\sqrt{2}$ in any experiment would falsify RS.
    \item Superluminal signaling via entanglement would falsify
    the $c = \ell_0/\tau_0$ speed limit.
    \item A local realistic theory reproducing all quantum
    correlations would undermine the shared-entry ontology.
  \end{itemize}
\end{enumerate}

%% ============================================================
\section{Conclusion}
\label{sec:conclusion}
%% ============================================================

Bell inequality violation follows from the shared-ledger
structure of entanglement.  The classical bound $|S| \leq 2$
fails because entangled particles share one global ledger
entry, not two local states.  The Tsirelson bound
$|S| \leq 2\sqrt{2}$ coincides with $\sqrt{2^D}$ for $D = 3$,
and the optimal measurement angle
$\pi/4 = 2\pi/2^D$ is one tick of the 8-tick epoch---a
structural coincidence whose deeper significance remains an
open question (see Remark~\ref{rem:d3status}).
No-signaling is a theorem of the $J$-cost phase invariance and
the ledger propagation speed.

%% ============================================================
\begin{thebibliography}{99}

\bibitem{RS_Axioms}
J.~Washburn, ``The Algebra of Reality: A Recognition Science Derivation
of Physical Law,'' \emph{Axioms} \textbf{15}(2), 90 (2026).
\texttt{doi:10.3390/axioms15020090}

\bibitem{Bell1964}
J.~S.~Bell, ``On the Einstein Podolsky Rosen Paradox,''
Physics \textbf{1}, 195 (1964).

\bibitem{CHSH1969}
J.~F.~Clauser, M.~A.~Horne, A.~Shimony, and R.~A.~Holt,
``Proposed Experiment to Test Local Hidden-Variable Theories,''
Phys.\ Rev.\ Lett.\ \textbf{23}, 880 (1969).

\bibitem{Tsirelson1980}
B.~S.~Tsirelson, ``Quantum Generalizations of Bell's
Inequality,'' Lett.\ Math.\ Phys.\ \textbf{4}, 93 (1980).

\bibitem{Aspect1982}
A.~Aspect, P.~Grangier, and G.~Roger,
``Experimental Realization of Einstein-Podolsky-Rosen-Bohm
Gedankenexperiment: A New Violation of Bell's Inequalities,''
Phys.\ Rev.\ Lett.\ \textbf{49}, 91 (1982).

\bibitem{Hensen2015}
B.~Hensen \textit{et al.}, ``Loophole-Free Bell Inequality
Violation Using Electron Spins Separated by 1.3~km,''
Nature \textbf{526}, 682 (2015).

\bibitem{Giustina2015}
M.~Giustina \textit{et al.}, ``Significant-Loophole-Free
Test of Bell's Theorem with Entangled Photons,''
Phys.\ Rev.\ Lett.\ \textbf{115}, 250401 (2015).

\bibitem{Shalm2015}
L.~K.~Shalm \textit{et al.}, ``Strong Loophole-Free Test
of Local Realism,'' Phys.\ Rev.\ Lett.\ \textbf{115},
250402 (2015).

\bibitem{Aczel1966}
J.~Acz\'el, \emph{Lectures on Functional Equations and Their
Applications}, Academic Press (1966).

\bibitem{PopescuRohrlich1994}
S.~Popescu and D.~Rohrlich, ``Quantum Nonlocality as an Axiom,''
Found.\ Phys.\ \textbf{24}, 379 (1994).

\end{thebibliography}

\end{document}
